\chapter{Abstract}
This work analyzes the main discrete-variable quantum key distribution protocols, with the aim of systematically characterizing their foundations, strengths, and weaknesses. A structured literature review methodology was applied, using a uniform template for each protocol covering its physical principles, operational scheme, known vulnerabilities, and practical examples. The protocols studied are: BB84, B92, E91, six-state protocol, SARG04, decoy-state BB84, LM05, and MDI-QKD. Results reveal a clear evolution from simple prepare-and-measure schemes toward entanglement-based and measurement-device-independent proposals, each addressing specific limitations of its predecessors. It is concluded that no single protocol is universally optimal: the choice depends on the implementation scenario, available technological capabilities, and the assumed threat model, with decoy-state protocols and MDI-QKD representing the most mature solutions for real-world deployments.

\vspace{0.5cm}
{\bf Keywords:} quantum key distribution, quantum cryptography, discrete-variable protocols, quantum entanglement, decoy states, measurement-device independence

\mainmatter